\section{Test: Factorització de polinomis}

\iniciTest{pol2}{a,b,c,b,d,a,c,b,d,c}

\begin{enumerate}

\pregunta Si un polinomi $p(x)$ compleix $p(a)=0$ per a un cert nombre real,
aleshores

\begin{enumerate}[label=\alph*)]

\opcio{a}{$p(x)$ és divisible per $x-a$}

\opcio{b}{$p(x)$ és divisible per $x+a$}

\opcio{c}{$p(x)$ no és divisible per $x-a$}

\opcio{d}{No es pot assegurar res si no sabem el valor de $a$}

\end{enumerate}

\feedback

\pregunta En la següent divisió\begin{equation*}
\begin{array}{c|c}
p(x) & x+a \\ \hline
R & c(x)\end{array}\end{equation*}es compleix:

\begin{enumerate}[label=\alph*)]

\opcio{a}{$p(a)=c(a)$}

\opcio{b}{$p(-a)=R$}

\opcio{c}{$p(a)=R$}

\opcio{d}{$p(-a)=c(a)$}

\end{enumerate}

\feedback

\pregunta Si un polinomi $p(x)$ té tots els seus coeficients enters i $-3$ és
una de les seves arrels, aleshores

\begin{enumerate}[label=\alph*)]

\opcio{a}{El seu terme independent pot ser $8$}

\opcio{b}{El seu terme independent pot ser $9$}

\opcio{c}{El seu terme independent pot ser qualsevol nombre negatiu}

\opcio{d}{$p(3)=0$}

\end{enumerate}

\feedback

\pregunta El quocient $c(x)$ i el residu $r(x)$ de dividir $\sqrt{2}x^{4}-4\sqrt{2}$ per $x+\sqrt{2}$ són

\begin{enumerate}[label=\alph*)]

\opcio{a}{$c(x)=\sqrt{2}x^{3}+2x^{3}+2\sqrt{2}x+4$ i $r(x)=0$}

\opcio{b}{$c(x)=\sqrt{2}x^{3}+2x^{3}+2\sqrt{2}x+4$ i $r(x)=\sqrt{2}$}

\opcio{c}{$c(x)=\sqrt{2}x^{3}-2x^{2}+2\sqrt{2}x-4$ i $r(x)=0$}

\opcio{d}{$c(x)=\sqrt{2}x^{3}-2x^{2}+2\sqrt{2}x-4$ i $r(x)=-\sqrt{2}$}

\end{enumerate}

\feedback

\pregunta Quin dels polinomis següents és divisible per $x-2$?

\begin{enumerate}[label=\alph*)]

\opcio{a}{$x^{2}+2x+1$}

\opcio{b}{$x^{2}+x-2$}

\opcio{c}{$x^{3}-3x+2$}

\opcio{d}{$x^{6}-64$}

\end{enumerate}

\feedback

\pregunta Si el residu de dividir $p(x)=x^{3}+ax^{2}-bx+6$ per $x-1$ és $5$ i el
de dividir-lo per $x+1$ és $3$, aleshores:

\begin{enumerate}[label=\alph*)]

\opcio{a}{$a=-2$ i $b=0$}

\opcio{b}{$a=-2$ i $b=-4$}

\opcio{c}{$a=2$ i $b=-4$}

\opcio{d}{$a=2$ i $b=-2$}

\end{enumerate}

\feedback

\pregunta Un polinomi té residu $5$ en dividir-lo per $x-3$ i també al
dividir-lo per $x+1$. Quin és el residu de dividir-lo per $(x-3)(x+1)$?

\begin{enumerate}[label=\alph*)]

\opcio{a}{$3x+5$}

\opcio{b}{$x+5$}

\opcio{c}{$5$}

\opcio{d}{No podem trobar-lo perquè falten dades}

\end{enumerate}

\feedback

\pregunta Les solucions reals de l'equació\begin{equation*}
x^{6}-x^{5}-6x^{4}-x^{2}+x+6=0 
\end{equation*}són:

\begin{enumerate}[label=\alph*)]

\opcio{a}{$-2$, $-1\,$, $1$ i $2$}

\opcio{b}{$-2$, $-1\,$, $1$ i $3$}

\opcio{c}{$-3$, $-1\,$, $1$ i $2$}

\opcio{d}{Cap de les anteriors}

\end{enumerate}

\feedback

\pregunta La descomposició factorial del polinomi $p(x)=4x^{4}-17x^{2}+4$ és

\begin{enumerate}[label=\alph*)]

\opcio{a}{$4(x-1)(x+1)(x-2)(x+2)$}

\opcio{b}{$4(x^{2}+1)(x-2)(x+2)$}

\opcio{c}{$4(x^{2}+1)(x^{2}+4)$}

\opcio{d}{$4(x-2)(x+2)(x-\frac{1}{2})(x+\frac{1}{2})$}

\end{enumerate}

\feedback

\pregunta Quin dels polinomis següents té $-2$ com a arrel doble, $0$
com a arrel triple i $1$ com a arrel de multiplicitat $4$?

\begin{enumerate}[label=\alph*)]

\opcio{a}{$(x-2)^{2}(x+1)^{4}$}

\opcio{b}{$x(x+2)^{2}(x-1)^{4}$}

\opcio{c}{$x^{3}(x+2)^{2}(x-1)^{4}$}

\opcio{d}{$x^{3}(x-2)^{2}(x+1)^{4}$}

\end{enumerate}

\feedback

\end{enumerate}

\bigskip
